\documentclass[11pt]{article}
\usepackage{fontspec,amsmath,amssymb,amsthm,hyperref}
\usepackage[margin=1in]{geometry}
\setlength{\emergencystretch}{3em}
\title{An intermediate bound with a supplied stationary Brownian reference}
\author{LangevinLean mathematical exposition}
\date{}
\begin{document}\maketitle
\noindent Source module: \texttt{UnadjustedLangevin/ULAFromBrownian.lean}. Compile this proof document with XeLaTeX or LuaLaTeX. The accompanying HTML places the complete Lean declarations beside every proof.


This auxiliary endpoint already derives the integral equation, continuity and moment bounds from the Brownian coordinates and potential assumptions. It keeps stationarity of the constructed reference segments as an explicit premise.


\section{Apply the reference estimate to the constructed segments}


The Brownian-reference constructor supplies stationary segments, score bound \(\sqrt{Ld}\), exact noise norm \(\sqrt{2d}\sqrt t\), and endpoint increments matching the actual Euler innovations. The potential comparison theorem derives \(m\le L\). Apply the small-step analytic reference theorem and simplify the start identity \(S_0(0)=Y_0\) to obtain

\[W_2(\operatorname{Law}(X_k),\pi)\le(1-mh)^k\|X_0-Y_0\|_2+\frac{33}{20}\frac Lm\sqrt{hd}.\]

This statement uses the supplied initial realization, not an infimum over couplings. Its stationarity premise is removed by CanonicalULA, where independent Brownian sources and the Gibbs-stationarity theorem are assembled.



\noindent Lean declarations: \texttt{dalalyan\_of\_constructed\_brownian\_reference}.

\end{document}
